\begin{frame}{자주 묻는 질문 (정리)}
  \begin{itemize}
    \item \kb{PCA·word2vec의 ``압축''과 같은가?} 방향은 비슷(고차원 원본을
      다루기 쉬운 형태로). PCA·word2vec은 구조 보존(비지도), DQN은
      ``Q값 정확히 예측''이라는 \textbf{지도} 목표.
    \item \kb{함수근사 = 비표 기반?} 사실상 같다 ---
      각 상태-행동을 \textbf{별개 파라미터}(표 칸)로 가지는 것 vs
      \textbf{파라미터를 공유하는 함수}로 표현. 선형 회귀·트리·RBF도
      비표 기반이고, DQN은 그중 \textbf{신경망}을 고른 것.
    \item \kb{일반화 = 과적합의 정반대?} 맞다. 너무 작은 망 =
      \textbf{부족 적합}, 너무 큰 망 = 배운 상태에만 과잉 적합 $\to$
      보지 못한 상태에서 오히려 더 나빠짐.
  \end{itemize}
\end{frame}
